% !PRINTED none
\weberpage{none}
% [source note: this Abschnitt opener has no printed folio; sequence position corresponds to inferred page 670.]

\begin{center}
\gesperrt{Twenty-Third Chapter.}

\gesperrt{Class Number.}

\rule{4em}{0.4pt}

\S.\,185.

\gesperrt{Dirichlet's Theorem on the Units.}
\end{center}

The investigations concerning the number of ideal classes in an algebraic field, which are now to give us the necessary completions of the proofs of the preceding chapter, make use of the methods that \gesperrt{Dirichlet} introduced into number theory\begingroup\renewcommand{\thefootnote}{\arabic{footnote})}\footnote{\gesperrt{Dirichlet's}, Recherches sur diverses applications de l'analyse infinitésimale à la théorie des nombres. \gesperrt{Crelle's} Journal, Bd. 19, 21. \gesperrt{Dirichlet's} Werke, Bd. I (1839, 1840). (Here belong, moreover, the posthumous memoirs of \gesperrt{Gauss} as well, \glqq De nexu inter multitudinem classium etc.\grqq{} \gesperrt{Gauss'} Werke, Bd. II.) The application of these principles to the cyclotomic numbers was first made by \gesperrt{Kummer} [\gesperrt{Crelle's} Journal, Bd. 35, 40 (1847, 1850); \gesperrt{Liouville's} Journal, Bd. 16 (1851)]. The general theory of the units was likewise founded by \gesperrt{Dirichlet} [\gesperrt{Dirichlet's} Werke, Bd. I, S. 622, 633, 639 (1840, 1842, 1846)]. The general theory of the units and the determination of the class numbers was given by \gesperrt{Dedekind}: \gesperrt{Dirichlet-Dedekind}, Vorlesungen über Zahlentheorie, 4. Auflage, \S.\,183 f. \gesperrt{Minkowski}, Geometrie der Zahlen.}\endgroup.

They are no longer purely algebraic, inasmuch as transcendental functions and numbers occur in them as well, such as the natural logarithm and the number $\pi$. Use is also made of passages to the limit, as in the integral calculus, from which we have already drawn a beautiful result in \S.\,156.

The investigations we have first to deal with are entirely general; that is, they hold for every algebraic number field, and it offers no advantage in simplicity to restrict them to special fields, such as the cyclotomic fields.
